\chapter{Priorisierung: Gewichtete Rankings und PageRank}
\label{chap:priorisierung}

\begin{myremark}[Algorithmusart in diesem Kapitel]
    Funktionsart: \textbf{Priorisierung}. \\
    Umsetzungsart: \textbf{regelbasiert und selbstlernend}.
\end{myremark}

\begin{myremark}
    Dieses Kapitel behandelt Algorithmen, die nicht primär Kategorien
    oder Zahlen vorhersagen, sondern Elemente in eine \textbf{Reihenfolge} bringen.
\end{myremark}

\section{Grundidee: Ordnen statt klassifizieren}

Bei einer Priorisierung wird nicht gefragt \enquote{Welche Klasse?} oder \enquote{Welcher Wert?},
sondern \enquote{Was soll zuerst kommen?}. Das Ergebnis ist eine \textbf{Rangliste}.

Typische Beispiele sind Suchmaschinen, Newsfeeds, Lernplattformen oder Support-Tickets.
Die Qualität eines Priorisierungsalgorithmus hängt davon ab, wie sinnvoll die Reihenfolge ist.

\section{Gewichtete Priorisierung}
Die gewichtete Priorisierung ist eine einfache Methode, um mehrere Kriterien zu einem Score zusammenzufassen. Dieser Algorithmus kommt zum Beispiel bei der Bewertung von Lernquellen (welches Buch, welcher Artikel, welches Video ist am besten geeignet?), bei der Auswahl von Support-Tickets (welches Problem soll zuerst gelöst werden?) oder bei der Priorisierung von Aufgaben (welche To-Do hat die höchste Dringlichkeit?) zum Einsatz.

\begin{mydefinition}[Gewichtete Priorisierung]
    Bei einer gewichteten Priorisierung bekommt jedes Kriterium ein \textbf{Gewicht}.
    Aus den Teilwerten $z_1, z_2, \ldots, z_n$ und Gewichten $w_1, w_2, \ldots, w_n$ entsteht ein Score:
    \[
        \text{Score}(x) = w_1 z_1 + w_2 z_2 + \cdots + w_n z_n = \sum_{i=1}^n w_i z_i
    \]
    Oft gilt $\sum_i w_i = 1$ und alle Teilwerte sind auf $[0,1]$ normiert.
\end{mydefinition}

\begin{mycaution}
    Wenn Merkmale unterschiedliche Skalen haben, verzerren grosse Zahlen die Rangliste.
    Darum müssen die Werte vor dem Zusammenfassen oft normalisiert werden.
    Für negative Kriterien wie \enquote{Aufwand} oder \enquote{Kosten} wird der Wert invertiert,
    z.\,B. durch $1-z$.
\end{mycaution}

\begin{myexample}
    Drei Quellen für ein Schulprojekt werden nach vier Kriterien bewertet:
    Relevanz, Aktualität, Vertrauenswürdigkeit und geringer Aufwand.
    Die Gewichte seien $0.40$, $0.25$, $0.25$ und $0.10$.
\end{myexample}

\begin{myexercise}[Gewichtete Scores berechnen]
    Drei Lernquellen wurden bereits auf $[0,1]$ normiert:

    \begin{center}
        \begin{tabular}{lcccc}
            \toprule
            Quelle & Relevanz & Aktualität & Vertrauen & Aufwand \\
            \midrule
            A      & 0.85     & 0.50       & 0.75      & 0.30    \\
            B      & 0.75     & 0.85       & 0.70      & 0.20    \\
            C      & 0.80     & 0.65       & 0.95      & 0.60    \\
            \bottomrule
        \end{tabular}
    \end{center}

    Berechnen Sie den Score mit
    \[
        \text{Score} = 0.40\cdot R + 0.25\cdot A + 0.25\cdot V + 0.10\cdot(1-\text{Aufwand})
    \]
    und ordnen Sie die drei Quellen.
    \begin{myanswer}
        A: $0.40\cdot0.85 + 0.25\cdot0.50 + 0.25\cdot0.75 + 0.10\cdot0.70 = 0.7225$ \\
        B: $0.40\cdot0.75 + 0.25\cdot0.85 + 0.25\cdot0.70 + 0.10\cdot0.80 = 0.7675$ \\
        C: $0.40\cdot0.80 + 0.25\cdot0.65 + 0.25\cdot0.95 + 0.10\cdot0.40 = 0.7600$ \\
        Rangfolge: B, C, A.
    \end{myanswer}
\end{myexercise}

\begin{myexercise}[Gewichte diskutieren]
    Was verändert sich, wenn man Aktualität stärker gewichtet als Vertrauenswürdigkeit?
    Nennen Sie ein Beispiel, in dem diese Entscheidung sinnvoll sein kann.
\end{myexercise}

\section{PageRank: Priorisierung im Netzwerk}

\begin{mydefinition}[PageRank]
    PageRank priorisiert Webseiten nicht nur nach der Anzahl eingehender Links,
    sondern auch nach der \textbf{Qualität} dieser Links.
    Eine Seite ist wichtig, wenn wichtige Seiten auf sie verweisen.
\end{mydefinition}

Die Grundidee ist ein \enquote{zufälliger Surfer}:
Mit Wahrscheinlichkeit $d$ folgt er einem Link, mit Wahrscheinlichkeit $1-d$ springt er zufällig auf eine neue Seite.
Der PageRank-Score einer Seite $p$ ist dann:
\[
    PR(p) = \frac{1-d}{N} + d \sum_{q \to p} \frac{PR(q)}{\text{outdeg}(q)}
\]
mit $N$ Seiten und Dämpfungsfaktor $d$ (typisch: $0.85$). Mit outdeg$(q)$ ist die Anzahl ausgehender Links von Seite $q$ gemeint. Mit $PR(q)$ ist der PageRank von Seite $q$ gemeint. Dabei wird $PR(q)$ anfänglich für alle Seiten gleich gesetzt (z.\,B. $1/N$) und iterativ aktualisiert, bis die Werte stabil sind.

\begin{mycaution}
    PageRank ist nicht neutral: Linkfarmen (also Seiten, die sich gegenseitig verlinken), gekaufte Links oder starke Netzwerke können die Rangfolge verzerren.
    Darum braucht man Missbrauchserkennung und zusätzliche Qualitätsmerkmale.
\end{mycaution}

\begin{myexercise}[PageRank intuitiv lesen]
    Betrachten Sie das Netzwerk:
    \begin{itemize}
        \item A verlinkt auf B und C
        \item B verlinkt auf C
        \item C verlinkt auf A
        \item D verlinkt auf C
    \end{itemize}
    Wenn alle Seiten mit gleichem PageRank starten und $d=0.85$ gilt:
    \begin{enumerate}
        \item Welche Seite dürfte nach einem Rechenschritt am höchsten priorisiert sein?
        \item Warum ist ein Link von einer Seite mit vielen Ausgängen weniger wert?
    \end{enumerate}
    \begin{myanswer}
        Nach einem Iterationsschritt ist C am höchsten priorisiert, weil A, B und D auf C verweisen.
        Ein Link von einer Seite mit vielen Ausgängen wird auf viele Ziele verteilt und zählt deshalb weniger stark.
    \end{myanswer}
\end{myexercise}

\begin{myexample}[Berechnung von PageRank]
    Um die PageRank-Werte für die Seiten A, B, C und D aus dem vorherigen Beispiel konkret zu berechnen, kann folgendermassen vorgegangen werden:


    Wir nehmen an, dass $d=0.85$ und setzen die Startwerte alle gleich auf $PR(A)=PR(B)=PR(C)=PR(D)=0.25$. Zunächst stellen wir die Linkstruktur als Matrix dar, wobei die Spalten die Zielseiten und die Zeilen die Quellseiten repräsentieren
    \[
        \begin{array}{c|cccc|c}
                          & A & B & C & D & \text{outdeg} \\
            \hline
            A             & 0 & 1 & 1 & 0 & 2             \\
            B             & 0 & 0 & 1 & 0 & 1             \\
            C             & 1 & 0 & 0 & 0 & 1             \\
            D             & 0 & 0 & 1 & 0 & 1             \\
            \hline
            \text{colsum} & 1 & 1 & 3 & 0 & 4             \\
        \end{array}
    \]
    Nun können wir die PageRank-Werte iterativ berechnen:
    \[
        \begin{aligned}
            PR(A) & = 0.0375 + 0.85 \cdot \left(\frac{PR(C)}{1}\right) = 0.0375 + 0.85 \cdot 0.25 = \underline{0.25}     \\
            PR(B) & = 0.0375 + 0.85 \cdot \left(\frac{PR(A)}{2}\right) = 0.0375 + 0.85 \cdot 0.125 = \underline{0.14375} \\
            PR(C) & = 0.0375 + 0.85 \cdot \left(\frac{PR(A)}{2} + \frac{PR(B)}{1} + \frac{PR(D)}{1}\right)               \\
                  & =0.0375 + 0.85 \cdot (0.125 + 0.25 + 0.25)                                                           \\
                  & =0.0375 + 0.85 \cdot 0.625 = \underline{0.56875}                                                     \\
            PR(D) & = 0.0375 + 0.85 \cdot 0 = \underline{0.0375}
        \end{aligned}
    \]
    Die Korrektheit der oben berechneten Werte kann überprüft werden, indem man die berechneten Werte summiert und sicherstellt, dass sie in etwa 1 ergeben (aufgrund von Rundungsfehlern kann es leicht abweichen). In diesem Fall ergibt die Summe $PR(A) + PR(B) + PR(C) + PR(D) = 0.25 + 0.14375 + 0.56875 + 0.0375 = 1.0$, was bestätigt, dass die Werte korrekt berechnet wurden.
\end{myexample}

\begin{myexercise}[PageRank selber berechnen]
    Sie haben in Ihrem Freundeskreis vier Personen: Alice, Bob, Carol und Dave. Ihre Freundschaften sind wie folgt:
    \begin{itemize}
        \item Alice ist mit Bob und Carol befreundet.
        \item Bob ist mit Alice und Dave befreundet.
        \item Carol ist mit Alice befreundet.
        \item Dave ist mit Bob befreundet.
    \end{itemize}
    Berechnen Sie die PageRank-Werte für jede Person, wenn $d=0.85$ gilt und alle Personen mit einem Startwert von $0.25$ beginnen.

    \begin{myanswer}
    Zunächst stellen wir die Freundschaftsstruktur als Matrix dar:
    \[
        \begin{array}{c|cccc|c}
                          & Alice & Bob & Carol & Dave & \text{outdeg} \\
            \hline
            Alice         & 0     & 1   & 1     & 0    & 2             \\
            Bob           & 1     & 0   & 0     & 1    & 2             \\
            Carol         & 1     & 0   & 0     & 0    & 1             \\
            Dave          & 0     & 1   & 0     & 0    & 1             \\
            \hline
            \text{colsum} & 2     & 2   & 1     & 1    & 6             \\
        \end{array}
    \]
    Nun berechnen wir die PageRank-Werte iterativ:
    \[
    \begin{aligned}
        PR(Alice) & = 0.0375 + 0.85 \cdot \left(\frac{0.25}{2} + \frac{0.25}{1}\right) = 0.0375 + 0.85 \cdot 0.375 = \underline{0.35625} \\
        PR(Bob)   & = 0.0375 + 0.85 \cdot \left(\frac{0.25}{2} + \frac{0.25}{1}\right) = 0.0375 + 0.85 \cdot 0.375 = \underline{0.35625} \\
        PR(Carol) & = 0.0375 + 0.85 \cdot \left(\frac{0.25}{2}\right) = 0.0375 + 0.85 \cdot 0.125 = \underline{0.14375} \\
        PR(Dave)  & = 0.0375 + 0.85 \cdot \left(\frac{0.25}{2}\right) = 0.0375 + 0.85 \cdot 0.125 = \underline{0.14375}
    \end{aligned}
    \]
    Nach einem Iterationsschritt haben Alice und Bob den höchsten PageRank-Wert von $0.35625$, gefolgt von Carol und Dave mit jeweils $0.14375$. Die PageRank-Werte summieren sich auf $0.35625 + 0.35625 + 0.14375 + 0.14375 = 1.0$.
    \end{myanswer}
\end{myexercise}

\begin{mychallenge}[Priorisierung vergleichen]
    Vergleichen Sie gewichtete Priorisierung und PageRank:
    \begin{enumerate}
        \item Welche Methode ist leichter erklärbar?
        \item Welche Methode ist anfälliger für Manipulation?
        \item Wo würden Sie welche Methode im Schulkontext einsetzen?
    \end{enumerate}
\end{mychallenge}

\begin{myremark}[Notebook]
    Praktische Umsetzung inkl. Übungen:
    \href{\notebookbaseurl06_Priorisierung_PageRank.ipynb}{\texttt{06\_Priorisierung\_PageRank.ipynb}}
\end{myremark}
